\hypertarget{advanced-text-processing-string-textwrap}{%
\chapter{\texorpdfstring{Advanced Text Processing (\texttt{string},
\texttt{textwrap})}{Advanced Text Processing (string, textwrap)}}\label{advanced-text-processing-string-textwrap}}

While \passthrough{\lstinline!str!} methods cover basic needs, the
\passthrough{\lstinline!string!} and \passthrough{\lstinline!textwrap!}
modules handle complex formatting and layout logic, often interacting
with terminal dimensions and internationalization.

\hypertarget{string.formatter-the-engine-of-.format}{%
\subsection{\texorpdfstring{43.1 \texttt{string.Formatter}: The Engine
of
\texttt{.format()}}{43.1 string.Formatter: The Engine of .format()}}\label{string.formatter-the-engine-of-.format}}

The \passthrough{\lstinline!f-string!} (Chapter 12) is the fastest, but
\passthrough{\lstinline!string.Formatter!} is the most extensible. *
\textbf{\passthrough{\lstinline!parse(format\_string)!}}: This method
returns an iterator of
\passthrough{\lstinline!(literal\_text, field\_name, format\_spec, conversion)!}.
* \textbf{\passthrough{\lstinline!get\_value(key, args, kwargs)!}}: This
is the hook for custom lookup logic. * \textbf{Internals}: F-strings are
compiled to specialized bytecode
(\passthrough{\lstinline!FORMAT\_VALUE!}), whereas
\passthrough{\lstinline!.format()!} calls into the
\passthrough{\lstinline!string!} module's C-accelerated formatting
logic.

\hypertarget{textwrap-dynamic-layout-management}{%
\subsection{\texorpdfstring{43.2 \texttt{textwrap}: Dynamic Layout
Management}{43.2 textwrap: Dynamic Layout Management}}\label{textwrap-dynamic-layout-management}}

\passthrough{\lstinline!textwrap!} is essential for CLI tools that must
adapt to varying terminal widths. *
\textbf{\passthrough{\lstinline!TextWrapper!} Object}: Maintains state
for \passthrough{\lstinline!width!}, \passthrough{\lstinline!indent!},
and \passthrough{\lstinline!break\_long\_words!}. *
\textbf{\passthrough{\lstinline!wrap()!}
vs.~\passthrough{\lstinline!fill()!}}: \passthrough{\lstinline!wrap!}
returns a list of strings; \passthrough{\lstinline!fill!} returns a
single newline-joined string. * \textbf{Algorithms}: It uses a greedy
algorithm to fit words into the specified width, handling edge cases
like hyphenated words and double-width Unicode characters correctly.

\begin{center}\rule{0.5\linewidth}{0.5pt}\end{center}
